\chapter{1978 Nobel Prize in Economics}
\textbf{Laureates: }Herbert Simon (USA)

\section*{Reason for Award}
Theory decision-making within economic organizations

\section{What They Did}
Herbert Simon introduced bounded rationality concept arguing human decision-making limited cognitive capacity information processing constraints satisficing behavior seeking satisfactory rather than optimal solutions organizational decision processes involve sequential search gradual improvement Organizational theory emphasized administrative behavior procedural rationality instead of perfect optimization

\section{Impact on Economic Thinking}
Simon fundamentally altered understanding rationality economics replacing perfect rationality bounded rationality creating foundation behavioral organizational studies decision research challenging traditional neoclassical assumptions His work showed how organizations make decisions under uncertainty information limitations emphasizing incrementalism satisficing heuristics These insights revealed how actual decision making differs from idealized models explaining organizational inefficiencies path dependence inertia

\section{Modern Perspectives Today}
Bounded rationality underlies behavioral economics prospect theory mental accounting heuristics biases Simon's organizational insights inform corporate governance strategy artificial intelligence cognitive science distributed problem-solving Decision-theoretic models incorporate cognitive limits heuristic strategies Behavioral organizational studies examine how groups make decisions Simon pioneered interdisciplinary approach combining psychology sociology economics studying organizations realistic way
